\documentclass[tikz,border=8pt]{standalone}
\usepackage{tikz}
\usepackage{amsmath}
\usepackage{pgfplots}
\pgfplotsset{compat=1.18}
\usepackage{ctex}
\usetikzlibrary{calc, positioning, arrows.meta, decorations.pathreplacing}

\begin{document}
\begin{tikzpicture}[scale=1.1]

  % 背景网格
  \draw[gray!20, thin] (-0.3,-0.3) grid (6.3,4.8);
  \draw[->,gray!50, thin] (-0.3,0) -- (6.5,0) node[right,font=\scriptsize] {$x$};
  \draw[->,gray!50, thin] (0,-0.3) -- (0,5.0) node[above,font=\scriptsize] {$y$};

  % 取 a1=(2,1), a2=(1,3), x=(1.5, 0.8)
  % Ax = 1.5*(2,1) + 0.8*(1,3) = (3+0.8, 1.5+2.4) = (3.8, 3.9)

  % 列向量 a1
  \draw[->,thick,blue!80!black,>=Stealth] (0,0) -- (2,1)
    node[below right,font=\small] {$\mathbf{a}_1=(2,1)$};
  % 列向量 a2
  \draw[->,thick,teal!80!black,>=Stealth] (0,0) -- (1,3)
    node[left,font=\small] {$\mathbf{a}_2=(1,3)$};

  % x1*a1 = 1.5*(2,1) = (3,1.5)
  \draw[->,thick,blue!50,>=Stealth,densely dashed] (0,0) -- (3,1.5)
    node[above,font=\scriptsize] {$x_1\mathbf{a}_1$};

  % x2*a2 = 0.8*(1,3) = (0.8,2.4)
  \draw[->,thick,teal!50,>=Stealth,densely dashed] (0,0) -- (0.8,2.4)
    node[left,font=\scriptsize] {$x_2\mathbf{a}_2$};

  % 从 x1*a1 出发，平移 x2*a2，到达 Ax
  \draw[->,densely dashed,orange!60,thick,>=Stealth] (3,1.5) -- (3.8,3.9);
  % 从 x2*a2 出发，平移 x1*a1，到达 Ax
  \draw[->,densely dashed,orange!60,thick,>=Stealth] (0.8,2.4) -- (3.8,3.9);

  % 结果向量 Ax
  \draw[->,very thick,red!80!black,>=Stealth] (0,0) -- (3.8,3.9)
    node[above right,font=\small] {$A\mathbf{x}=(3.8,\,3.9)$};

  % 子空间平面（a1 和 a2 张成的线，因为是 R^2 所以整个平面，用阴影带暗示）
  % 画出 a1, a2 张成的方向线（延伸线，浅色）
  \draw[blue!20, thin] (-0.3,-0.15) -- (5.3,2.65);  % a1 方向
  \draw[teal!20, thin] (-0.15,-0.45) -- (1.65,4.95); % a2 方向

  % 原点
  \fill (0,0) circle (1.5pt) node[below left,font=\scriptsize] {$O$};
  \fill (3.8,3.9) circle (2pt);

  % 说明框（右下角，远离所有向量）
  \node[font=\small, align=center, fill=white, draw=gray!50, thin,
        rounded corners=2pt, inner sep=6pt]
    at (5.5, 0.9)
    {$A\mathbf{x} = x_1\mathbf{a}_1 + x_2\mathbf{a}_2$\\
     永远在 $\mathbf{a}_1,\mathbf{a}_2$\\张成的子空间中};

  % 标题
  \node[font=\large\bfseries] at (3.0,-1.0) {矩阵乘法的列空间视角};

\end{tikzpicture}
\end{document}
